\section{An adapter must preserve the scientific object}
\label{exec14:l6-adapter}

Typed input is useful only if the adapter preserves what the fields mean.
Consider two implementations of the same analysis path: a direct local
service and a framework-mediated service. Let $x$ be a valid request with
its scientific context, and let $F(x)$ be the result of a qualified
deterministic engine under that context. An adapter encodes the request
with $E$, sends it through the tool interface, and decodes the response
with $D$. The desired comparison is semantic:
\begin{equation}
 D\{F_{\rm tool}(E(x))\}\equiv F(x),
 \label{exec14:l6-commute}
\end{equation}
where $\equiv$ requires the same target, assumptions, qualification and
result meaning, not identical incidental ordering of JSON fields. For a
pure deterministic numerical path it may also require exact values, or a
declared numerical tolerance supported by the computation. For stochastic
methods the comparison must declare randomization and distributional
criteria. Passing one exact fixture does not prove this relation for every
request. This is a proposed qualification relation with explicit domains,
not an asserted property of an installed framework.
\citep{exec14_l6_adapter_exact}

Several superficially convenient encodings cannot satisfy that relation.
Suppose an encoder replaces every finite or unbounded effect set by a pair
containing its smallest and largest endpoints. Then the union
$(-\infty,-2]\cup[3,\infty)$ and the whole line both produce the same
pair of infinite endpoints. A downstream decoder receiving only that pair
cannot know which input was supplied. If it returns the union, it is wrong
for the whole-line input; if it returns the whole line, it is wrong for the
union as an exact representation. It could return a conservative outer
approximation if the contract explicitly permits loss of precision, but
it must label that transformation rather than claim equivalence. An exact
finite representation of endpoint pairs reproduces this collision in the
project example.

There is a general reason. Let $E(x_1)=E(x_2)$ for two scientifically
distinguishable objects $x_1$ and $x_2$. Any deterministic decoder $D$ has
the same output on the two equal encodings. Therefore both identities
$D(E(x_1))=x_1$ and $D(E(x_2))=x_2$ cannot hold. Randomized decoding cannot
guarantee recovery of both without additional information either. The
same argument applies when an adapter collapses PAG endpoint types, maps
unknown and nonidentified into one status, or drops which population an
effect targets. Before considering a port, the proposal should inspect
whether the receiving schema is expressive enough. More elaborate retry
logic cannot reconstruct information the schema discarded.
\citep{exec14_l6_adapter_exact}

A refusal state also carries scientific information. A graphical engine
can return an identification expression, a formal nonidentification result
within a specified model class, an unsupported method, or an operational
failure. These states call for different continuation actions. A framework
that treats every returned string as a successful tool result may preserve
transport bytes while discarding that distinction in its next prompt.
The adapter should expose a tagged result and require the next operation
to consume an admitted tag. A report can state that a method is unsupported
without claiming the target is unknowable. A timeout cannot become an
estimated effect of zero. These conditions can be tested with authored
fixtures before a model is selected.

\subsection{Units, signs and set-valued banking decisions}

Unit conversion illustrates a transformation that can preserve meaning
when it is explicit. Let an effect $\theta$ be measured in dollars and
suppose an interface returns thousands of dollars: $\theta'=\theta/1000$.
An interval or union must be transformed as a set, not only by changing
a display label. More generally, for a nonzero constant $a$ and constant
$b$, let $T(\theta)=a\theta+b$. For a closed interval $[l,u]$, its image
is $[al+b,au+b]$ when $a>0$ and $[au+b,al+b]$ when $a<0$. Infinite
endpoints and each component of a union follow the same ordering rule.
Empty and whole-line sets retain their distinct identities.
\citep{exec14_l6_adapter_exact}

If a set $C$ has coverage at least $1-\alpha$ for $\theta$, then its exact
image $T(C)$ has coverage at least $1-\alpha$ for $T(\theta)$ when $T$ is
a fixed known bijection: the events $\theta\in C$ and
$T(\theta)\in T(C)$ are identical. An outer enclosure of $T(C)$ can
preserve the lower coverage guarantee while losing precision. A numerical
inner approximation has no such automatic guarantee. If the conversion
rate or cost is estimated, the fixed-known-map premise fails and its
uncertainty must be accounted for. This elementary argument explains why
representation and numerical error belong in an adapter contract. It
does not validate the original set's sampling law.

Take a stipulated offer-minus-no-offer effect set
$C=(-\infty,-2]\cup[3,\infty)$ dollars and an incremental offer cost
of one dollar. The net-value set is
$C-1=(-\infty,-3]\cup[2,\infty)$. Its worst case is negative and its
best case is positive, so this set alone does not justify either robust
acceptance or robust rejection under a sign rule. Reporting only the
midpoint of a finite substitute interval would invent a decision. If a
receiving tool instead uses the reverse treatment contrast, the effect
set becomes $-C=(-\infty,-3]\cup[2,\infty)$, but this numerical
coincidence in the chosen example does not make reversing the contrast
equivalent to subtracting cost. The object must retain the target and
transformation, not only the endpoints.

For a less coincidental bounded check, let $C=[2,5]$ and a positive
one-dollar cost. Net value is $[1,4]$, so every admitted value supports
the offer under the stipulated expected-value sign criterion. Reversing
the effect contrast produces $[-5,-2]$; it is not $[1,4]$. Scaling
dollars into thousands maps the original effect to $[.002,.005]$ and
the cost to $.001$ in those units. Subtracting a cost of one in the new
units would be a thousandfold error. Exact rational images in the example
check the sign, units and component ordering. The resulting conclusions
remain illustrative consequences of stipulated sets and costs, not
validated lending advice or empirical financial benefit.
\citep{exec14_l6_adapter_exact}

\subsection{Porting needs a failure-domain comparison}

A port's acceptance domain should include both supported and deliberately
unsupported requests. In supported cases compare scientific context,
endpoint conventions, uncertainty tags, set topology, estimates and
applicable numerical tolerances. In unsupported cases compare the reason
and retained provenance. Test NaN, infinity, missing timestamps, duplicate
variable identifiers, a changed sign convention and a graph with unresolved
endpoints according to the receiving method's contract. The expected
answer may be refusal. A test that merely requires the port to return
something for every case rewards the very semantic substitutions the
product is meant to prevent.

The comparison should record which differences are representational and
which are substantive. Object-key order may be irrelevant to the engine
yet consequential to a byte digest; a normalized serialization can be
introduced with an explicit version and migration. Changing a dependency
law from independent accounts to monthly panels is substantive even if
the numerical array is unchanged. A port can preserve syntax and still
change inference. These distinctions motivate a thin typed adapter first,
followed by a port only when measured operational benefits justify the
additional qualification burden.

\section{Attribute improvements to the component that changed}
\label{exec14:l6-factorial}

Comparing an elaborate new agent to an old script can conflate at least
three changes: a stronger backend, better orchestration and access to
better causal tools. A product decision needs to know which component
accounts for an observed improvement and at what cost. A prospective
comparison can hold documents, permitted tools, review rubric and total
repair opportunity fixed while varying backend and orchestration. Let
$b\in\{A,B\}$ index two candidate backends and $f\in\{0,1\}$ a plain
controller versus the framework candidate. For eligible task $i$, define
loss $L_i(b,f)$ to include the predeclared scientific error and operational
failure consequences. Lower loss is better.

If all four outcomes can be obtained without contamination, a paired
task comparison isolates the framework contrast at each backend. If
repeated exposure changes human judgment or leaks a reference answer,
use independent adjudicators, blinded outputs or randomized assignment
with a declared estimand. Neither approach permits tuning the framework
on the final task answers. Tasks that share a generating mechanism are
not automatically independent observations, so task families and their
sampling law must be declared before attaching inferential uncertainty.
The prospective construction here is an evaluation design, not a claim
of collected experimental outcomes.
\citep{exec14_l6_adapter_exact}

Write $\mu_{bf}=E{L_i(b,f)}$ for the mean under the target evaluation
task distribution. The mean framework contrast, averaging equally over
the two backends, is
\begin{equation}
 \Delta_F=\tfrac12[(\mu_{A1}-\mu_{A0})+
                              (\mu_{B1}-\mu_{B0})],
\end{equation}
and the analogous backend contrast is
\begin{equation}
 \Delta_B=\tfrac12[(\mu_{B0}-\mu_{A0})+
                              (\mu_{B1}-\mu_{A1})].
\end{equation}
The interaction
$\Delta_{BF}=(\mu_{B1}-\mu_{B0})-(\mu_{A1}-\mu_{A0})$
measures how the framework contrast changes with backend. These are
descriptive contrasts of the stipulated factorial task experiment. Their
identification from sampled trials depends on the assignment and missing
outcome rules. Unequal deployment frequencies would require the relevant
weights rather than the equal weights used here.

An exact example takes mean losses
$(\mu_{A0},\mu_{A1},\mu_{B0},\mu_{B1})=(10,8,7,6)$.
Then $\Delta_F=-1.5$, $\Delta_B=-2.5$ and $\Delta_{BF}=1$.
The framework reduces mean loss by two units with A but only one with B.
The diagonal comparison $\mu_{B1}-\mu_{A0}=-4$ mixes both changes;
it cannot all be attributed to orchestration. The average contrasts
sum to the diagonal difference, yet their separate magnitudes carry the
investment information. Costs require their own four-cell measurement,
including failed tasks and review effort. The numbers are authored
illustrations, not performance claims about any named framework.
\citep{exec14_l6_adapter_exact}

Availability creates a further denominator problem. Suppose both routes
receive 100 tasks. Route A returns 90 answers with 81 correct; route B
returns 60 with 57 correct. B's correctness among returned answers is
95 percent, above A's 90 percent, but its correct-answer fraction among
all assigned tasks is 57 percent versus 81 percent. Neither summary
alone settles utility when correct refusals and harmful wrong answers
have different values. If every unanswered task incurs a stipulated
loss of three and every wrong returned answer a loss of ten, A incurs
120 total loss units and B 150. At loss one per unanswered task, the
totals become 100 and 70. The preferred route changes with the actual
decision consequences. These exact counts show why missingness and
refusal must be classified before ranking an architecture.

A properly specified causal task sometimes requires a refusal or a set
rather than a scalar estimate. Such an output belongs in the correct
answer category if the rubric calls for it. Backend refusal for an
unrelated service-policy reason, inability to parse, and unavailable
scientific method are different categories. They should remain visible
even if a decision maker later aggregates their losses. The prospective
pilot should first establish that the rubric recognizes these outcomes,
then compare useful assistance under that rubric. A subsequent bank
trial can ask whether the resulting architecture preserves meaning and
improves decisions in its own population; the public/generated study
cannot silently stand in for that external question.
\citep{exec14_l6_adapter_exact}

\label{exec14:l6-adapter-end}
